\section{Methodology}
\label{sec:methodology}

\subsection{Overview}

Building on the theoretical intuition that spurious-minority samples should have gradients conflicting with the majority-dominated batch mean, we define a \textbf{Gradient Alignment Score} as the negated cosine similarity between a sample's per-sample last-layer gradient and the within-batch mean gradient.
A low alignment score (conflicting direction) serves as the minority membership signal --- the core of the proposed Gradient Alignment Debiasing (GAD) framework.

We measure whether the gradient alignment score achieves higher ROC-AUC than per-sample loss as a binary predictor of spurious-minority group membership, at multiple training epochs.
The intervention (upweighting) is not evaluated here --- it is only motivated if the directional signal surpasses loss.

\subsection{Per-Sample Gradient Alignment Score}

Let $f_\theta: \mathcal{X} \to \mathbb{R}^C$ be a neural network.
For a mini-batch $\mathcal{B} = \{(x_i, y_i)\}_{i=1}^{B}$, define per-sample last-layer gradients:
\begin{equation}
  g_i = \nabla_{\theta_{\text{fc}}} \mathcal{L}(f_\theta(x_i), y_i)
\end{equation}
and batch-mean gradient $\bar{g} = \frac{1}{B} \sum_{i=1}^{B} g_i$.
The alignment score is:
\begin{equation}
  a_i = 1 - \frac{g_i \cdot \bar{g}}{\|g_i\|_2 \cdot \|\bar{g}\|_2} \in [0, 2]
\end{equation}
where high $a_i$ indicates gradient conflict with the batch mean (predicted minority signal).
ROC-AUC of $a_i$ below 0.5 means minority samples have \emph{lower} conflict scores --- equivalently, higher raw cosine similarity with the batch mean --- the opposite of the directional hypothesis.

\textbf{Rationale for last-layer scope:} Restricting to $\theta_{\text{fc}}$ (Linear(2048, $C$), 4,098 parameters for $C=2$) provides computational feasibility with vmap and directly captures class-discriminative gradient direction.

\textbf{Rationale for cosine similarity:} Magnitude-invariant, isolating the directional component from per-sample loss magnitude --- the theorized advantage over loss-based proxies.

\subsection{Efficient Computation via vmap}

\begin{verbatim}
from torch.func import functional_call, vmap, grad

def per_sample_loss(params, buffers, x, y):
    pred = functional_call(model, (params, buffers), (x.unsqueeze(0),))
    return F.cross_entropy(pred, y.unsqueeze(0))

ft_per_sample_grad = vmap(grad(per_sample_loss), in_dims=(None, None, 0, 0))
fc_params = {k: v for k, v in model.named_parameters() if 'fc' in k}
per_sample_grads = ft_per_sample_grad(fc_params, {}, x_batch, y_batch)
\end{verbatim}

The vmap scope is restricted to \texttt{model.fc}; the backbone is frozen during probe computation.
Per-batch gradient storage: $32 \times 4{,}098 = 131\text{K}$ floats.

\subsection{Probe Injection Protocol}

At each checkpoint epoch $e \in \{1, 5, 10, 25, 50\}$, after standard ERM training:
\begin{enumerate}
  \item Freeze model parameters.
  \item Run full pass over training set (batch size 32).
  \item Compute $a_i$ and $\ell_i$ for all $N$ training samples.
  \item Evaluate ROC-AUC of each against binary minority group membership labels.
\end{enumerate}

\subsection{Training Configuration}

\begin{table}[h]
\centering
\caption{Training hyperparameters.}
\label{tab:training}
\begin{tabular}{lll}
\toprule
\textbf{Parameter} & \textbf{Waterbirds} & \textbf{CelebA} \\
\midrule
Model & ResNet-50 (IMAGENET1K\_V1) & ResNet-50 (IMAGENET1K\_V1) \\
Optimizer & SGD, lr=0.001, mom=0.9, wd=1e-4 & SGD, lr=0.0001, mom=0.9, wd=1e-4 \\
Batch size & 32 & 32 \\
Train size & 4,795 & 16,000 (stratified subsample) \\
Seed & 42 & 42 \\
\bottomrule
\end{tabular}
\end{table}
